\section{Introduction}
\label{sec:introduction}

Pre-trained Large Language Models (LLMs), such as GPT~\cite{achiam2023gpt} and LLaMa~\cite{llama}, have demonstrated strong capabilities across diverse tasks, including text generation~\cite{brown2020gpt3,zhao2023survey}, question answering~\cite{lewis2020retrieval, guu2020retrieval}, and code generation~\cite{chen2021evaluating-code,nijkamp2022codegen}. 
To adapt these models to personalized or domain-specific tasks, \textit{fine-tuning} on pre-trained LLMs is typically performed.
Such adaptation is essential for scenarios like biomedical analysis~\cite{zhang2024generalist, tian2024opportunities}, personalized chatbot interactions~\cite{wang2022self-instruct}, or specialized customer support~\cite{yun2023fine-tuning-dialogue-summarization}. 
However, traditional full-model fine-tuning, where all model parameters are learned, requires substantial hardware resources such as multiple nodes, each node equipped with multiple flagship GPUs (e.g., NVIDIA B200 GPUs~\cite{B200}) and interconnected by high-speed links (e.g., NVLink~\cite{nvlink} and InfiniBand~\cite{infiniband}).  
For instance, full-model fine-tuning of LLaMa-3.1-70B~\cite{llama3.1} requires approximately 1120GB of GPU memory for model states alone (parameters, gradients, optimizer states), making this approach prohibitively expensive for practical applications~\cite{llama3.1-training-cost}.

\begin{sloppypar}
To mitigate the substantial hardware requirements, Parameter-Efficient Fine-Tuning (PEFT) methods~\cite{hu2022lora,li2021prefix-tuning,lester2021prompt-tuning,kopiczko2023vera,liu2024dora,dettmers2023qlora,li2023loftq,ding2023sparse-low-rank-adaptation,houlsby2019peft-learning-nlp,lin2020exploring,ruckle2020adapterdrop,zhao2024apt,zhang2025memory-efficient-lora} have emerged. These approaches significantly reduce resource usage by keeping pre-trained LLM parameters \textit{frozen} (not updated during fine-tuning) and selectively updating only a small set of injected trainable parameters, known as \textit{adapters}.
Among these methods, Low-Rank Adaptation (LoRA)~\cite{hu2022lora} and its variants~\cite{liu2024dora,dettmers2023qlora,li2023loftq,ding2023sparse-low-rank-adaptation} are particularly popular due to their simplicity and effectiveness. 
%
% LoRA freezes the model's pre-trained weights and introduces two small, low-rank linear layers as residual adapters, fine-tuning only these adapters.
% Typically, LoRA adapters contain fewer than 1\% of the original model parameters, substantially lowering memory requirements for storing gradients and optimizer states. 
LoRA freezes the model's pre-trained weights and adds a residual branch composed of two trainable low-rank linear layers: a down-projection from the input dimension $k$ to a smaller rank $r$, followed by an up-projection back to the output dimension $n$, where $r \ll min(n, k)$. 
This low-rank branch is added to the output of the original frozen layer, enabling task-specific adaptation without modifying the base model (see Figure~\ref{fig:background-lora-illustration} and Section~\ref{subsec:background-llm-lora-finetuning}).
%
For example, fine-tuning LLaMa-3.1-70B~\cite{llama3.1} using LoRA adapters with a low-rank dimension of 16 introduces only 0.29\% additional parameters, reducing GPU memory usage to 142GB while preserving strong model quality~\cite{dettmers2023qlora,zhang2025memory-efficient-lora}.
\end{sloppypar}

Due to these substantial resource savings, LoRA fine-tuning has become widely adopted on both cloud platforms~\cite{openai-fine-tuning,together-ai-fine-tuning,runpod-fine-tuning,fireworks-fine-tuning} and local environments~\cite{zheng2024llamafactory,llama-cookbook}.
To realistically apply LoRA in practice and fine-tune high-quality adapters, users often run multiple jobs in parallel. These jobs may explore different hyperparameter settings~\cite{tribes2023hyperparameter-instruction-tuning} or continuously adapt models to evolving datasets and user needs~\cite{jindal2024balancing,ibrahim2024simple-continual-pre-training}.
As a result, fine-tuning throughput, measured as the number of training samples processed per second, has become a key metric for reducing both cost and overall training time.

However, despite significant algorithmic advances in LoRA-based approaches~\cite{dettmers2023qlora,kopiczko2023vera,liu2024dora} and explorations in serving scenarios~\cite{zhou2022pets,chen2024punica,sheng2023s-lora,wu2024dlora}, existing LoRA fine-tuning systems for LLMs largely reuse optimizations designed for full-model fine-tuning.
Specifically, systems like PEFT Library~\cite{peft}, LLaMA-Factory~\cite{zheng2024llamafactory}, and llama-cookbook~\cite{llama-cookbook} typically rely on a subset of parallelization methods such as Fully Sharded Data Parallelism (FSDP)~\cite{pytorch-fsdp,zero}, Tensor Parallelism (TP)~\cite{megatron-1,flexflow}, or Pipeline Parallelism (PP)~\cite{gpipe,pipedream,terapipe} to fit the training into GPU memory and achieve efficiency with multiple GPUs.
While such techniques are still useful, our analysis reveals they do not sufficiently address the unique characteristics of LoRA fine-tuning, causing significant inefficiencies.

The first limitation is the high runtime overhead introduced by LoRA adapters. 
LoRA adapters add less than 1\% parameters compared to the full model and are thus expected to incur minimal computation and memory overhead. 
However, our profiling shows that applying LoRA reduces training throughput by approximately 40\% compared to the original frozen model. 
This overhead comes from increased memory traffic: the small LoRA projection layers are memory-bandwidth-bound due to their low-rank dimensions, and the operations in the adapter repeatedly load and store large activation tensors. 
These factors together increase global memory access by up to $2.64\times$, as detailed in Section~\ref{sec:motivation-lora-overhead}.

The second limitation is the lack of efficient support for joint fine-tuning across multiple LoRA adapters. 
Typically, each adapter is fine-tuned independently, even if they share the same base model. 
Since LoRA adapters are lightweight and add minimal memory footprint pressure, multiple jobs can be combined and run jointly on the same GPUs. 
We refer to this strategy as \textit{multi-LoRA fine-tuning}~\cite{ye2023mlora}, which is highly practical for hyperparameter tuning~\cite{tribes2023hyperparameter-instruction-tuning} and multi-tenant cloud services~\cite{fireworks-fine-tuning,together-ai-fine-tuning}.
Although multi-LoRA optimization~\cite{chen2024punica,sheng2023s-lora,wu2024dlora} has been widely used in LLM serving, the motivation in our setting is entirely different. 
Serving systems batch requests to increase arithmetic intensity during autoregressive single-token decoding~\cite{chen2024punica,sheng2023s-lora}, whereas fine-tuning already processes full sequences with sufficient arithmetic intensity.
Instead, the primary benefit of multi-LoRA fine-tuning comes from mitigating the overhead of distributed parallelism~\cite{ye2023mlora}.

\begin{sloppypar}
While existing multi-LoRA fine-tuning systems like mLoRA~\cite{ye2023mlora} can reduce pipeline bubbles by filling them with independent groups of samples from different adapters, we observe that this approach is incomplete and still suffers from significant inefficiencies.
First, it still relies on generic LoRA kernels that are bottlenecked by redundant memory accesses to large activation tensors as discussed previously, and batching more samples does not help.
Second, it fails to address the load imbalance across GPUs.
Realistic fine-tuning workloads often contain samples with variable sequence lengths, leading to imbalanced work across GPUs.
This imbalance creates idle time in data-parallel replicas and increases pipeline bubbles from poorly aligned microbatches.
As we analyze in Section~\ref{sec:motivation-multi-lora-oversight}, strategically grouping and scheduling samples across multiple jobs is critical to mitigate this imbalance and unlock further efficiency gains.
\end{sloppypar}

Based on these insights, we argue that an efficient LoRA fine-tuning system must both reduce the runtime overhead of LoRA adapters and leverage multi-LoRA optimization opportunities to reduce distributed training overhead.
To meet these requirements, we propose \textit{\Name}, an efficient multi-level fusion system tailored specifically for LoRA fine-tuning of LLMs.
%
At the kernel level, our key insight is that most of LoRA's overhead comes from memory-bandwidth-bound operations on large activation tensors.
We address this by splitting the computation graph at the point where the tensor size shrinks to the low-rank dimension $r$.
This allows us to fuse memory-bandwidth-bound operations without recomputation or synchronization, while preserving the optimal performance of compute-bound matrix multiplications. 
This design leads to our \texttt{FusedLoRA} and \texttt{FusedMultiLoRA} kernels.
%
At the scheduling level, \Name enables concurrent fine-tuning of multiple LoRA adapters that share the same base model by batching samples across jobs.
Compared to the existing multi-LoRA systems, this further improves system throughput by reducing pipeline bubbles and balancing GPU workloads.
The challenge is that such batching must preserve execution dependencies between global batches, especially under pipeline parallelism.
To address this, we first group adapters in a way that creates natural gaps between their batches, then solve a bin-packing problem within each group to construct balanced, dependency-safe microbatches using a combination of Mixed Integer Linear Programming (MILP) and greedy heuristics.

We implement \Name on top of Megatron-LM~\cite{megatron-1}, a state-of-the-art distributed training framework, to support efficient parallelism and scalability.
We then extensively evaluate \Name on a wide range of LLMs, including LLaMa-3.1-8B, Qwen-2.5-32B, and LLaMa-3.1-70B, across NVIDIA H100 and L40S GPUs with various realistic datasets. 
Results show that \Name achieves up to $1.96\times$ ($1.47\times$ on average) end-to-end speedup compared to Megatron-LM with FSDP and PP, and up to $1.46\times$ ($1.29\times$ on average) improvement over mLoRA~\cite{ye2023mlora}, the state-of-the-art multi-LoRA fine-tuning system.
Additionally, our fused kernel alone achieves up to $1.39\times$ ($1.27\times$ on average) speedup compared to the default LoRA implementation on NVIDIA H100 GPUs, and can directly serve as a plug-in replacement in existing systems, offering immediate benefits to the broader community.

Overall, this paper makes the following contributions:

\begin{figure}[t]
    \centering
    \includegraphics[width=0.99\linewidth]{figures/background-fig-1-lora.pdf}
    \hspace{0.01\linewidth} (a) Full Fine-tuning
    \hspace{0.125\linewidth} (b) LoRA Fine-tuning
    \hspace{0.125\linewidth}
    \caption{Comparison between traditional full model fine-tuning and LoRA fine-tuning.}
    \Description{Comparison between traditional full model fine-tuning and LoRA fine-tuning. Traditional fine-tuning updates all model parameters, while LoRA fine-tuning only updates lightweight adapter matrices.}
    \label{fig:background-lora-illustration}
\end{figure}

\begin{sloppypar}
\begin{itemize}
    \item We identify two key limitations in existing LoRA fine-tuning systems: high runtime overhead from redundant memory access, and missed opportunities for optimizing multi-job training. To address both, we propose \Name, a multi-level fusion system for accelerating LLM fine-tuning on modern GPU clusters.
    \item We propose a horizontal fusion strategy that reduces redundant memory access without disrupting compute-bound performance. We also design a scheduling algorithm that groups adapters across fine-tuning jobs and batches their samples to improve GPU load balance and reduce pipeline overhead.
    \item We evaluate \Name across diverse LLMs, datasets, and GPU platforms, showing significant improvements over existing LoRA fine-tuning systems. Our fused kernel also provides strong standalone gains and can be directly integrated into more general LoRA systems.
\end{itemize}
\end{sloppypar}
